% Options for packages loaded elsewhere
\PassOptionsToPackage{unicode}{hyperref}
\PassOptionsToPackage{hyphens}{url}
\documentclass[
]{article}
\usepackage{xcolor}
\usepackage[margin=1in]{geometry}
\usepackage{amsmath,amssymb}
\setcounter{secnumdepth}{-\maxdimen} % remove section numbering
\usepackage{iftex}
\ifPDFTeX
  \usepackage[T1]{fontenc}
  \usepackage[utf8]{inputenc}
  \usepackage{textcomp} % provide euro and other symbols
\else % if luatex or xetex
  \usepackage{unicode-math} % this also loads fontspec
  \defaultfontfeatures{Scale=MatchLowercase}
  \defaultfontfeatures[\rmfamily]{Ligatures=TeX,Scale=1}
\fi
\usepackage{lmodern}
\ifPDFTeX\else
  % xetex/luatex font selection
  \setmainfont[Ligatures=TeX]{Cambria}
  \setmonofont[]{Consolas}
\fi
% Use upquote if available, for straight quotes in verbatim environments
\IfFileExists{upquote.sty}{\usepackage{upquote}}{}
\IfFileExists{microtype.sty}{% use microtype if available
  \usepackage[]{microtype}
  \UseMicrotypeSet[protrusion]{basicmath} % disable protrusion for tt fonts
}{}
\makeatletter
\@ifundefined{KOMAClassName}{% if non-KOMA class
  \IfFileExists{parskip.sty}{%
    \usepackage{parskip}
  }{% else
    \setlength{\parindent}{0pt}
    \setlength{\parskip}{6pt plus 2pt minus 1pt}}
}{% if KOMA class
  \KOMAoptions{parskip=half}}
\makeatother
\usepackage{longtable,booktabs,array}
\newcounter{none} % for unnumbered tables
\usepackage{calc} % for calculating minipage widths
% Correct order of tables after \paragraph or \subparagraph
\usepackage{etoolbox}
\makeatletter
\patchcmd\longtable{\par}{\if@noskipsec\mbox{}\fi\par}{}{}
\makeatother
% Allow footnotes in longtable head/foot
\IfFileExists{footnotehyper.sty}{\usepackage{footnotehyper}}{\usepackage{footnote}}
\makesavenoteenv{longtable}
\usepackage{graphicx}
\makeatletter
\newsavebox\pandoc@box
\newcommand*\pandocbounded[1]{% scales image to fit in text height/width
  \sbox\pandoc@box{#1}%
  \Gscale@div\@tempa{\textheight}{\dimexpr\ht\pandoc@box+\dp\pandoc@box\relax}%
  \Gscale@div\@tempb{\linewidth}{\wd\pandoc@box}%
  \ifdim\@tempb\p@<\@tempa\p@\let\@tempa\@tempb\fi% select the smaller of both
  \ifdim\@tempa\p@<\p@\scalebox{\@tempa}{\usebox\pandoc@box}%
  \else\usebox{\pandoc@box}%
  \fi%
}
% Set default figure placement to htbp
\def\fps@figure{htbp}
\makeatother
\setlength{\emergencystretch}{3em} % prevent overfull lines
\providecommand{\tightlist}{%
  \setlength{\itemsep}{0pt}\setlength{\parskip}{0pt}}
\usepackage{xurl}
\usepackage{bookmark}
\IfFileExists{xurl.sty}{\usepackage{xurl}}{} % add URL line breaks if available
\urlstyle{same}
\hypersetup{
  pdftitle={Frontier and Localhost: How Production AI Learns Outside the Weights (Short Version)},
  pdfauthor={Mikhail L. Arbuzov; Alexey A. Shvets; Sisong Bei},
  hidelinks,
  pdfcreator={LaTeX via pandoc}}

\title{Frontier and Localhost: How Production AI Learns Outside the
Weights (Short Version)}
\author{Mikhail L. Arbuzov \and Alexey A. Shvets \and Sisong Bei}
\date{}

\begin{document}
\maketitle

\textbf{Abstract}

Production LLM systems increasingly adapt outside the model weights.
After deployment failures, teams modify prompts, rules, memories,
skills, tools, eval suites, routing graphs, and governance pipelines on
a cadence that frontier weight updates cannot match. But this scaffold
layer is still mostly maintained as \emph{patchwork}: fixes are proposed
by intuition, committed with weak credit assignment, accumulated without
pruning, and promoted beyond the scope where they were validated. This
paper formalises a disciplined alternative, which we call
\emph{artifact-layer descent}. The pipeline is concrete. Recurring
residuals are assigned to scaffold coordinates. Candidate artifact
deltas are tested against patch loss. Accepted deltas persist with
rollback. Promotion across contexts is bounded by \emph{evidence
radius}: a local fix spreads only as far as the evidence supports.
Surveying approximately 130 production and research systems from 2023 to
2026, we find current systems partially instantiate this loop but leave
a central architecture gap, namely governed scaffold optimization across
contexts. The contribution is not the claim that prompts are weights. It
is to name the missing optimizer for an adaptation layer the field has
already built.

\begin{center}\rule{0.5\linewidth}{0.5pt}\end{center}

\subsection{1. The gradient moved outside the
model}\label{the-gradient-moved-outside-the-model}

The standard mental model of an LLM system is the model. Pre-training
fixes the weights, fine-tuning nudges them, deployment wraps a thin
prompt-plus-tool layer around the result and ships. Production in 2025
and 2026 has stopped behaving this way. Frontier models update on cycles
measured in months. The systems built on top of them update
continuously, in markdown files, persistent memories, retrieved skills,
registered tools, evaluation suites, agent topologies,
version-controlled prompt repositories, PR-gated governance. A growing
share of deployment-time adaptation now occurs outside the weights, in a
fast-changing \emph{local scaffold} wrapped around a slow-changing
frontier. The section title is a metaphor. The gradient inside the model
still happens during training, and the scaffold layer is not a gradient
in the calculus sense. The substantive claim is the milder one.
Production AI has accidentally created an external learning surface, and
right now it is being operated as \emph{patchwork} rather than as a
disciplined optimization process.

Read together, these artifacts form a single external adaptation surface
that current production practice still treats as scattered engineering.
Naming the surface makes the optimization discipline it needs (and
visibly lacks) visible. A Claude Skill is a learned procedural weight
kept outside the model. A Cursor Rule is a local parameter for one
project's distribution shift. A RAG connector is patch-specific memory
access. A tool registry is capability provisioning along an axis the
base model cannot reach. A PR review against a shared rule repository is
a promotion gate on a federated update. An eval suite that blocks merge
is a validation-loss check before promotion. A cross-tenant artifact
repository is \emph{federated artifact promotion} across distributed
shards: scope expansion under a gate, not arithmetic averaging (we
sharpen this in §3 and §9). The mechanisms were already in production.
What is new is reading them as a single external adaptation surface
rather than as scattered engineering choices, and asking what
optimization discipline that surface needs.

This paper surveys approximately 130 production and research systems,
identifies the six substrates where the scaffold lives, formalises the
two-loop architecture (local update, cross-context promotion under a
gate), and names a missing architecture as the most consequential open
problem: governed cross-tenant scaffold optimization with versioned
lineage, role-based access, and rollback. The single-line thesis is
this: \textbf{the field has built an external adaptation layer for LLMs,
but not the optimizer that should govern it.} The rest of the paper
develops that observation into a survey of where production sits on a
\emph{patchwork to random-search to descent} gradient, a formalisation
of the disciplined limit case under stated assumptions, and a named
missing architecture where the optimizer most plainly does not yet
exist.

\subsection{2. Why patch errors live outside the
weights}\label{why-patch-errors-live-outside-the-weights}

The scaffold is not merely where production systems happen to store
patches. For high-churn, tenant-specific, auditable adaptation, the
scaffold is \emph{often the practical substrate}. The reason is not that
no alternative exists. It is that under deployment-time constraints
(release cadence, tenant isolation, inspectability, rollback), the
scaffold is cheap, locally scoped, reversible, and updatable at human
time-scales while the weights are not.

\textbf{Frontier training optimises for breadth.} The weights are tuned
for cross-patch generalisation. They smooth over local conventions,
contradictory tenant preferences, and rare workflow-specific failures in
order to preserve broad competence across millions of unseen
deployments. The smoothing is the training objective, not an artefact.
RLHF reduces output diversity and homogenises preferred styles across
users (Kirk et al., 2024; Padmakumar \& He, 2023), and
instruction-tuning data composition rewards patterns that generalise
across the annotator pool rather than patterns specific to any one
deployment. The frontier model is therefore \emph{trained to be unable}
to encode a single tenant's idiosyncrasies as preferred behaviour.

\textbf{Production reliability is achieved in depth.} Inside a single
deployment patch, reliability is determined by exactly what global
training had to smooth away: local naming conventions, team-specific
workflows, idiosyncratic tool chains, customer-specific risk tolerances,
recurring local mistakes, hidden evaluator expectations. The per-patch
quantities from our previous work (failure-mode catalogue size,
hard-fraction, mode-discovery rate) are local by construction, and a
library calibrated to one patch under-covers the next.

\textbf{Encoding every patch in the weights is structurally infeasible.}
Three constraints rule out the obvious approach. \emph{Economic:}
millions of patches multiplied by billions of parameters times per-patch
fine-tuning cost is not a viable training-compute budget.
\emph{Release-cycle:} the patch wants daily updates; weight releases
ship quarterly at best. \emph{Cross-patch interference:} overfitting the
weights to one tenant's preferences degrades adjacent tenants whose
preferences contradict.

Per-tenant fine-tunes, LoRA adapters, distilled small models, and
learned routing each absorb some patch residuals in principle. In
practice the scaffold dominates, because it is the substrate that
satisfies all four production constraints at once: cheap, inspectable,
reversible, locally scoped, and updatable on a daily cadence rather than
a release cadence. The placement question that \emph{Architecture of
Errors} left open (where do the interventions live?) has this as its
answer. The frontier model generalises. The local scaffold specialises.
Reliability comes from governing that specialisation.

\textbf{Three stages of scaffold maintenance recur across the corpus.}
They differ in how candidate updates are produced, whether credit
assignment is explicit, and whether promotion is evidence-bounded. Most
production deployments sit at Stage 1 or 2. The math of §3 applies to
Stage 3.

{\def\LTcaptype{none} % do not increment counter
\begin{longtable}[]{@{}
  >{\raggedright\arraybackslash}p{(\linewidth - 8\tabcolsep) * \real{0.2000}}
  >{\raggedright\arraybackslash}p{(\linewidth - 8\tabcolsep) * \real{0.2000}}
  >{\raggedright\arraybackslash}p{(\linewidth - 8\tabcolsep) * \real{0.2000}}
  >{\raggedright\arraybackslash}p{(\linewidth - 8\tabcolsep) * \real{0.2000}}
  >{\raggedright\arraybackslash}p{(\linewidth - 8\tabcolsep) * \real{0.2000}}@{}}
\toprule\noalign{}
\begin{minipage}[b]{\linewidth}\raggedright
Stage
\end{minipage} & \begin{minipage}[b]{\linewidth}\raggedright
Update mechanism
\end{minipage} & \begin{minipage}[b]{\linewidth}\raggedright
Credit assignment
\end{minipage} & \begin{minipage}[b]{\linewidth}\raggedright
Promotion discipline
\end{minipage} & \begin{minipage}[b]{\linewidth}\raggedright
Dominant failure mode
\end{minipage} \\
\midrule\noalign{}
\endhead
\bottomrule\noalign{}
\endlastfoot
\textbf{1. Patchwork} & Intuition-driven local edits & Implicit / ad hoc
& None & Bloat \\
\textbf{2. Random artifact search} & Many variants tried under eval &
Weak; eval is the only signal & Eval-gated CI but no
failure-to-coordinate tracing & Overfitting to in-distribution patch \\
\textbf{3. Artifact-layer descent} & Residual to coordinate to delta &
Explicit per failure mode & Evidence-radius-bounded; reversible & (the
discipline this paper formalises) \\
\end{longtable}
}

The paper's central claim is not that modern AI systems are doing
stochastic descent. They mostly are not. The claim is that they
\emph{become so} once the patchwork loop is instrumented with measurable
patch loss, residual capture, failure-to-coordinate credit assignment,
candidate delta synthesis, and a gate that estimates directional loss
change before persistence or promotion. The survey of §§5--8 asks which
deployed systems have closed which of those five steps. §9 names the
closure that no public system has yet achieved.

\subsection{3. Artifact-layer descent}\label{artifact-layer-descent}

\textbf{The disciplined version of scaffold maintenance is a concrete
pipeline.} Observe a failure. Diagnose its failure mode. Identify which
scaffold coordinate caused it: a stale memory entry, an under-specified
instruction, a missing tool argument, a brittle retrieval rule, a
misrouted agent edge. Propose a candidate artifact edit. Test the edit
against patch loss. Accept or reject by a gate. Persist accepted edits;
prune or roll back when they regress. Promote across contexts only as
far as the evidence supports. We call that last constraint the edit's
\emph{evidence radius}: a local fix should only spread to other
deployments to the extent it has been validated there. The rest of this
section formalises this pipeline. §§4--9 ask which deployed systems have
closed which of its steps.

\textbf{Most current production does not close this loop.} A failure
happens, someone adds a prompt or rule or memory entry, the artifact
pile grows, nobody knows what caused improvement, nobody knows when to
prune, and fixes leak into broader contexts. This is patchwork (Stage 1
of §2), not stochastic descent. Even systems with eval-gated CI (Stage
2) usually lack failure-to-coordinate credit assignment, so improvements
are real but unattributed and brittle to transfer. Scaffold updates
\emph{become} stochastic descent (Stage 3) only when the patchwork loop
is instrumented with five things: measurable patch loss, residual
capture, failure-to-coordinate credit assignment, candidate delta
synthesis, and a gate that estimates directional loss change before
persistence or promotion. The descent picture below is mechanism when
those five conditions hold, and modeling language otherwise.

\textbf{Two loops, in plain English first.} \emph{Loop 1} is the local
fix: a failure produces a candidate edit, the edit is tested, the
accepted edit changes one deployment's scaffold. \emph{Loop 2} is the
decision whether that local fix should be shared with other deployments.
It promotes from one repository to a team-wide template, from one
tenant's playbook to a cross-tenant skill library, and so on. Loop 1
changes one \(L_D\). Loop 2 changes \emph{which} \(L_D\) is being
descended on.

\textbf{Assumptions and scope of the formalization.} The descent
argument below holds when: (i) the frontier model parameters
\(\theta_M\) are fixed during the adaptation window; (ii) scaffold edits
are persistent, inspectable artifacts distinct from transient context;
(iii) the patch admits a measurable operational loss; (iv) the proposal
process can attribute failures to scaffold coordinates with non-zero
probability; (v) gates are imperfect estimators of finite directional
loss change rather than oracle access to the true loss change; (vi)
Loop-2 promotion is scope expansion under a gate, not arithmetic
averaging of comparable updates. Survey claims of the form ``no surveyed
system has X'' are bounded by the audited public corpus and the survey
cutoff of May 2026.

\textbf{Patch scope hierarchy.} \emph{Patch} is used throughout this
paper with a defined scope hierarchy: USER ⊂ PROJECT ⊂ STACK ⊂ TENANT ⊂
CORE. A USER patch is one human's account or fork. A PROJECT is one
repository, workflow, or product workspace. A STACK is a curated bundle
of projects under one team. A TENANT is an organisation or customer.
CORE is the cross-tenant universal scope. Each level induces its own
loss as an average over the patches it contains. When the text says
``cross-tenant promotion'' it specifically means moving artifacts up to
the TENANT or CORE level. ``Patch-specific adaptation'' without
qualifier means descent on a USER or PROJECT-level loss.

The patch loss is the expected residual failure rate of the
model-scaffold composite on patch \(D\):

\[L_D(\theta_s) = \mathbb{E}_{x \sim D}\!\left[\ell(\pi(\theta_M, \theta_s, x))\right],\]

where \(\theta_s\) is a structured scaffold parameter object over six
substrate coordinates and \(\ell\) is an operational loss (failure
indicator, user-correction rate, eval error, schema violation, reviewer
disagreement). The important distinction is substrate. \(\theta_M\) is
fixed during deployment. \(\theta_s\) is the object being fit.

A failure does not yield an analytic gradient. It yields a candidate
edit. The gate \(G_t\) accepts the edit if its noisy estimate of the
directional loss change clears a margin. The accepted edit applies to
the current scaffold; rejected edits leave the scaffold unchanged:

\[\theta_s^{t+1} = \begin{cases} \theta_s^t \oplus z_t, & G_t(z_t)=1, \\ \theta_s^t, & G_t(z_t)=0, \end{cases}\]

with the gate accept condition

\[G_t(z_t)=1 \iff \widehat{\Delta}_{z_t} L_D(\theta_s^t) + \lambda\,\Omega(z_t) \leq -\tau_t,\]

where \(\widehat{\Delta}\) is the gate's estimator of the loss change,
\(\Omega \geq 0\) is a complexity penalty on the edit (instruction
length, tool-permission breadth, promotion scope), \(\lambda \geq 0\)
regularises that penalty, and \(\tau_t \geq 0\) is the required
improvement margin. Production systems implement this without writing
the equation: an eval blocks a merge on regression, a reviewer rejects a
vague rule, a governance system prevents a local artifact from being
shared until evidence accumulates. Different surveyed gate families
instantiate the abstraction differently. Eval-cascade, RL-threshold,
LLM-judge, surrogate verifier, and reviewer-judgment families directly
populate the per-edit gate. Semantic-merge and MAP-Elites archive
selection act as proposal-pruning steps that \emph{precede} the per-edit
gate. The per-system mapping is in Appendix C.

\textbf{Two loops, two distinct knobs.} Consider a concrete case. An
engineer notices that a coding agent keeps misnaming a class in one
repository. The fix is a one-line rule in that repo's instruction file.
The edit is small. Its scope is one project. The same engineer could
promote the rule to a team-wide template, or to the company's CORE
instructions, or to a cross-tenant skill registry. The \emph{edit} did
not change. The \emph{scope} did. And the gate that should let it
through gets stricter at each step, because the loss it is being tested
against is now an average over more patches, and a fix that helps one
repo can hurt another.

That observation maps cleanly onto two distinct quantities in the
optimisation analogy. \emph{Edit magnitude} (how much a single accepted
edit changes the local scaffold) plays the role of learning rate.
\emph{Promotion radius} (how many downstream patches the accepted edit
affects) controls \emph{which loss} is being descended on. USER-level
acceptance descends on one patch's \(L_D\). PROJECT promotion descends
on an average over a project's patches. TENANT or CORE promotion
descends on a federated average over many tenants' losses. The two knobs
are independently controllable, and the gate must be calibrated against
both. This is why governance is not administrative overhead. It is the
regularization system that prevents a small fix from being descended on
the wrong loss.

A note on Loop 2 and FedAvg. Loop 2 is \emph{federated artifact
promotion}: selection, merge, curation, review, distribution. It is not
arithmetic averaging of numeric updates. FedAvg is a useful analogy only
in the limit where artifacts have been embedded as mergeable update
representations and the aggregation is parameter-space averaging. No
surveyed system meets that limit.

Each loop has a gate. Gates are \textbf{Auto-gated} when the decision is
an operationalised algorithmic criterion (RL threshold, eval score,
LLM-judge with quantified agreement, surrogate verifier, benchmark
cascade). They are \textbf{Human-gated} when the decision is reviewer
judgment. They are \textbf{None} when the loop is absent. These five
terms (Loop 1, Loop 2, Auto-gated, Human-gated, None) anchor §§5--9.

The descent argument is \emph{mechanism} when five conditions hold:
measurable patch loss, residual capture, failure-to-coordinate credit
assignment, candidate delta synthesis, and a gate that estimates
directional loss change before persistence or promotion. It is
\emph{metaphor} when those conditions are absent. Appendix E formalises
the assumptions A1--A5 needed for a one-step descent inequality and a
convergence proposition of the loss to a \(G\)-stable scaffold, and
works through the RAG, pitch-review, tool-use, and memory cases that
make the mechanism visible.

\subsection{4. Six scaffold substrates}\label{six-scaffold-substrates}

A patch-local scaffold composes six substrates. They are the coordinates
of \(\theta_s\), each independently editable.

{\def\LTcaptype{none} % do not increment counter
\begin{longtable}[]{@{}
  >{\raggedright\arraybackslash}p{(\linewidth - 4\tabcolsep) * \real{0.3333}}
  >{\raggedright\arraybackslash}p{(\linewidth - 4\tabcolsep) * \real{0.3333}}
  >{\raggedright\arraybackslash}p{(\linewidth - 4\tabcolsep) * \real{0.3333}}@{}}
\toprule\noalign{}
\begin{minipage}[b]{\linewidth}\raggedright
Substrate
\end{minipage} & \begin{minipage}[b]{\linewidth}\raggedright
What persists
\end{minipage} & \begin{minipage}[b]{\linewidth}\raggedright
Score-5 mark
\end{minipage} \\
\midrule\noalign{}
\endhead
\bottomrule\noalign{}
\endlastfoot
\textbf{S1. Instructions} & Project/team/org rules attached to a
workspace & Versioned, two-loop, Auto-gated promotion \\
\textbf{S2. Skills} & Procedural artifacts (code/text) loaded on demand
& Skill bank self-modifies on failure, unit-test gated, versioned,
promoted \\
\textbf{S3. Memory} & Cross-session experience (episodic / semantic /
procedural) & Provenance + versioning + correction pathways \\
\textbf{S4. Tools} & Vetted action surface + context schemas &
Domain-curated bundle is the product; eval-gated tool versions \\
\textbf{S5. Orchestration} & Multi-role graph + routing &
Population/policy updates from outer-loop signal \\
\textbf{S6. Governance} & Versioning, promotion, audit, rollback &
dev→staging→prod with eval thresholds + audit log + rollback \\
\end{longtable}
}

The 0--5 rubric is uniform across substrates: \textbf{0} ephemeral;
\textbf{1} persistent local; \textbf{2} persistent + tool attachment, no
feedback; \textbf{3} multi-scope or feedback but no automated promotion;
\textbf{4} automated feedback updates the scaffold (one loop closed);
\textbf{5} two-loop versioned promotion plus governance (both loops
closed). We adopt this six-coordinate decomposition because each
coordinate is independently editable in production systems; alternative
partitions (e.g., separating evals from governance) are possible, and
the descent argument of §3 is invariant to the partition so long as the
independent-editability property holds.

\subsection{5. Survey method}\label{survey-method}

The survey covers approximately 130 unique LLM systems published or
productionised between January 2023 and May 2026. Of these,
approximately 90 satisfy the patch-local discriminator (score ≥ 3 on at
least one substrate) and form the \textbf{core corpus}. The remaining
\textasciitilde38 are commonly described as agentic or self-improving
but fail at least one discriminator; they are held aside as a
\textbf{contrast corpus}. A separate, stricter audit applied to 22
candidate systems uses a \emph{composite two-loop} criterion: Loop 1
modifies a persisted artifact from session signal \emph{and} Loop 2 has
an explicit cross-context promotion mechanism with versioning or
lineage. Thirteen research systems pass the strict audit. Two industry
exemplars (Sierra Agent OS 2.0, Cognition Devin) are also evaluated
using vendor documentation as primary evidence and flagged in §8. Full
inclusion/exclusion criteria, evidence-tier definitions (T1
peer-reviewed; T2 arXiv; T3 production claim; T4 open-source artifact;
T5 industry blog), scoring discipline, sensitivity analysis under two
natural relaxations of the criterion, and the representative per-system
audit trail are in Appendix B. Appendix A lists 30 representative survey
rows, and Appendix D lists representative contrast entries.

\subsection{6. Substrate maturity is
uneven}\label{substrate-maturity-is-uneven}

Aggregating the survey by substrate gives a maturity gradient: tools and
integrated systems are the most production-mature; skills and
orchestration lag; instructions are ubiquitous but shallow; memory is
widespread but rarely versioned. The unevenness is the survey's central
empirical pattern. The field has independently converged on the six
substrates, but no system matures all six at once. Research systems
learn fast and lack enterprise governance; production systems govern
well and learn slowly; open standards solve artifact portability but not
adaptation; vertical-bundle vendors solve patch specificity but rarely
expose cross-tenant promotion. Per-substrate score distributions,
score-5 exemplars, and the five surprising absences that fall out of the
discriminator (failure-triggered memory rare in production;
cross-customer skill promotion absent; memory versioning mostly
unsolved; scaffold-level curriculum unspecified; Loop-2 overfitting
detection missing) are in Appendix F. Each absence reappears as a
Paper-4 target in §11.

\subsection{7. Four clusters}\label{four-clusters}

Four archetypal clusters recur across the corpus. \textbf{Research
full-auto} systems automate both loops under algorithmic gates
(NanoResearch, SkillRL, AlphaEvolve, DGM, ADAS, Voyager). They
demonstrate that the architecture works; enterprise governance is
uniformly absent. \textbf{Production hybrid} systems pair an algorithmic
Loop-1 gate with a reviewer-judgment Loop-2 gate (SkillForge, CASCADE,
Sierra Agent OS 2.0, Cognition Devin) and represent the
production-viable cell. \textbf{Governance-first} systems automate Loop
2 only: eval-gated CI blocks promotion on regression, but Loop 1 is
absent because the candidate artifacts are engineer-authored
(Braintrust, Vellum, LangSmith Hub, W\&B Weave, Agenta, LangFuse). We
frame this as \emph{MLOps for scaffolds}, distinct from MLOps for
weights. \textbf{Open standards} systems converge across vendors on
artifact format and tool attachment (AGENTS.md in over 60,000 GitHub
repositories with measured efficiency gains; MCP as the cross-vendor
tool standard; Claude Skills, Cursor Rules, Copilot custom instructions,
Windsurf, Continue.dev, Zed, Replit on a common markdown-plus-tool
pattern). Convergence across vendors is the strongest evidence the
architecture is real rather than a single vendor's local optimum.
Per-cluster system walks with metrics, the orchestration-bounding result
of Tran \& Kiela (2026), and the vertical-bundle vendor placement are in
Appendix F.

\subsection{8. The two-loop design
space}\label{the-two-loop-design-space}

The composite-two-loop systems sit in a 3×3 matrix indexed by gate type
on each loop.

{\def\LTcaptype{none} % do not increment counter
\begin{longtable}[]{@{}
  >{\raggedright\arraybackslash}p{(\linewidth - 6\tabcolsep) * \real{0.2500}}
  >{\raggedright\arraybackslash}p{(\linewidth - 6\tabcolsep) * \real{0.2500}}
  >{\raggedright\arraybackslash}p{(\linewidth - 6\tabcolsep) * \real{0.2500}}
  >{\raggedright\arraybackslash}p{(\linewidth - 6\tabcolsep) * \real{0.2500}}@{}}
\toprule\noalign{}
\begin{minipage}[b]{\linewidth}\raggedright
\end{minipage} & \begin{minipage}[b]{\linewidth}\raggedright
\textbf{Loop-2 Auto-gated}
\end{minipage} & \begin{minipage}[b]{\linewidth}\raggedright
\textbf{Loop-2 Human-gated}
\end{minipage} & \begin{minipage}[b]{\linewidth}\raggedright
\textbf{Loop-2 None}
\end{minipage} \\
\midrule\noalign{}
\endhead
\bottomrule\noalign{}
\endlastfoot
\textbf{Loop-1 Auto-gated} & NanoResearch, SkillRL, AlphaEvolve, ADAS,
DGM, Voyager, SAGE, EvolveR, CoEvoSkills, AutoAgent, AutoManual &
SkillForge, CASCADE, Sierra OS 2.0†, Cognition Devin† & MAE*, MetaGen*,
MetaReflection*, CLIN* \\
\textbf{Loop-1 Human-gated} & \textbf{{[}empty by engineering logic{]}}
& \textbf{{[}empty in surveyed corpus{]}} & \emph{(none)} \\
\textbf{Loop-1 None} & ExpeL (weak), Braintrust*, LangSmith Hub*,
Vellum* & Anthropic skills repo*, DSPy*, TextGrad* & Self-Refine*,
ReAct*, Mem0* \\
\end{longtable}
}

\emph{Asterisks fail at least one composite-two-loop criterion; included
for contrast.} † Industry exemplar evaluated using vendor documentation
as primary evidence.

Three cells are populated. \textbf{{[}Auto × Auto{]}} holds eleven
research systems with machine-evaluated decision criteria on both loops.
\textbf{{[}Auto × Human{]}} is the production-viable cell: Loop 1 fires
at machine speed under an algorithmic criterion, Loop 2 requires human
approval before propagation. \textbf{{[}None × Auto{]}} holds the
governance-first prompts-as-code systems with eval-gated CI but no
failure-triggered candidate generation. Three cells are empty for
distinct structural reasons. \textbf{{[}Human × Auto{]}} is empty by
engineering logic: if humans author candidates manually, automated
promotion offers little leverage. \textbf{{[}Human × Human{]}} is empty
in the surveyed corpus; a candidate occupant would be a solo or
small-team workflow where the marginal cost of a bad auto-merge exceeds
the marginal cost of human latency, and no public example was found in
the May 2026 window. \textbf{{[}Human × None{]}} is structurally
improbable: a system requiring human authorship on Loop 1 without any
Loop 2 collapses to engineer-edits-a-config, and is not patch-local in
the sense of §3.

Two cross-cutting findings sharpen the picture. \emph{Noise reduction
via batching is largely absent in the surveyed auto-gated systems}.
Every auto-gated system fires at \(n = 1\) per generation step,
absorbing estimator noise statistically via the per-edit gate rather
than aggregating evidence across multiple proposals. Multi-sample
agreement filters (accept only when \(k\) independent proposals
converge) are a tractable design dimension that no surveyed system has
formalised. Their noise-reduction behaviour is qualitatively distinct
from mini-batch averaging (false-accept rate shrinks roughly as \(p^k\)
rather than variance shrinking as \(1/k\)), and the two should not be
conflated. \emph{Multi-layer hierarchies are under-explored.} Only
AlphaEvolve demonstrates explicit depth in the auto-gated cluster (via
cascade evaluator and MAP-Elites archive); the other twelve
composite-two-loop systems treat their artifact pool as flat. The full
cell-by-cell discussion, the closest-approaches enumeration that
supports §9, and the DGM dual-mode footnote are in Appendix F.

\subsection{9. The missing architecture: governed cross-tenant scaffold
optimization}\label{the-missing-architecture-governed-cross-tenant-scaffold-optimization}

Four properties define the architecture that no surveyed system fully
instantiates:

\begin{itemize}
\tightlist
\item
  \begin{enumerate}
  \def\labelenumi{(\alph{enumi})}
  \tightlist
  \item
    \textbf{Auto-gated Loop 1.} The commit criterion is an algorithmic
    threshold.
  \end{enumerate}
\item
  \begin{enumerate}
  \def\labelenumi{(\alph{enumi})}
  \setcounter{enumi}{1}
  \tightlist
  \item
    \textbf{Multi-tenant cross-org promotion.} Loop 2 routes artifacts
    between organisations.
  \end{enumerate}
\item
  \begin{enumerate}
  \def\labelenumi{(\alph{enumi})}
  \setcounter{enumi}{2}
  \tightlist
  \item
    \textbf{Versioned lineage with RBAC.} Promoted artifacts carry
    semantic versions, audit trails, and access-control policies.
  \end{enumerate}
\item
  \begin{enumerate}
  \def\labelenumi{(\alph{enumi})}
  \setcounter{enumi}{3}
  \tightlist
  \item
    \textbf{Rollback.} A bad promotion can be reverted without redeploy.
  \end{enumerate}
\end{itemize}

The closest approaches partition the requirements. AlphaEvolve has (a)
and partial (b) within a single organisation. SkillForge has (a) and (b)
within one enterprise with partial (c) via VFS commits. Sierra Agent OS
2.0 has (b) and partial (a). The governance-first cluster has (c) and
(d) explicitly but no (a) or (b) in the strict sense. The full
enumeration is in Appendix F.

In ML terms, what is missing is a privacy-preserving cross-tenant
aggregator with an automated eval gate at the scaffold layer. The
architecture would carry per-tenant artifact pools with tenant-tagged
provenance, aggregation across tenants with operator-tunable privacy
bounds, an automated eval gate blocking promotion when patch loss does
not decrease, versioned lineage with RBAC, and a rollback envelope on
every promoted artifact. The DP-FedAvg framing from the
federated-learning literature is a useful analogy for this target, but
only if artifacts are first embedded as mergeable update
representations. In surveyed practice, Loop 2 is selection and review
rather than parameter averaging.

Naming the missing architecture converts a vague gap into a
four-property audit criterion that any future system can be measured
against. Whether it \emph{should} be filled is deployment-dependent. In
safety-critical domains, reviewer-judgment gates may be a design feature
rather than a defect. But no public system currently fills it, and that
is what matters for the next two years of architecture work. Whether the
optimizer-equipped composite reaches deployment-level reliability is a
deployment-by-deployment empirical question; the paper's contribution is
to name what \emph{would have to be true} of the optimizer for that
question to even be testable.

\subsection{10. A new failure surface}\label{a-new-failure-surface}

Patchwork scaffold maintenance introduces failure modes weight-only
systems do not have. Each one is also a missing piece of the optimizer.

\emph{Scaffold bloat} is the dominant near-term risk. IFScale
(Jaroslawicz et al., 2025) reports latency growing roughly 35× for one
reasoning model (o4-mini) as instruction count scales from 10 to 250,
and the best frontier model in that benchmark (Gemini 2.5 Pro) drops to
68\% accuracy at 500 instructions, with weaker models degrading further.
The underlying mechanism (lost-in-the-middle, instruction interference;
Liu et al., 2023) persists across model generations even when individual
frontier-model scaling improves. The architectural mitigation is a
complexity penalty \(\Omega(z)\) in the gate and explicit pruning in the
loop.

\emph{Poisoned memory} is the dominant security risk. PoisonedRAG
reports 90\%+ attack success rate against standard RAG, and Memory
Control Flow Attacks report \textgreater90\% vulnerability across major
LLM agent frameworks. The architectural mitigation is provenance and
gating on memory writes, with rollback envelopes on accepted entries.

\emph{Prompt injection via tools and MCP} is the dominant cross-vendor
attack surface. CVE-2025-54136 (``MCPoison'') demonstrated a public
MCP-server compromise pathway. The architectural mitigation is
permission scoping and promotion gates that treat tool bindings as
first-class scaffold coordinates rather than ambient configuration.

\emph{Eval fragility} is the dominant ecosystem risk. The Leaderboard
Illusion documents up to 112\% relative performance gains under
differential training-data access. Eval suites are themselves
patch-local artifacts and must survive Goodhart's Law. The architectural
mitigation is held-out evals scoped within evidence radius rather than
promoted as universal benchmarks.

\emph{Coordination failures} are documented in MAST: 14 failure modes,
\(\kappa = 0.88\) agreement, dominantly role-confusion and context-loss.
The architectural mitigation is treating orchestration topology as a
scaffold coordinate subject to the same gate and rollback discipline as
instructions and memories.

\emph{Patch overfitting} is structural. Scaffolds tuned to a deployment
over-fit it. The architectural mitigation is \emph{evidence radius}
discipline at promotion: do not propagate a fix beyond the scope that
validated it.

\emph{Plasticity loss at the scaffold level} is the slow-moving risk.
Context rot, instruction conflict, memory pollution, and tool-version
drift produce a scaffold-level analogue of weight-level plasticity loss.
No surveyed system treats pruning or expiration as a first-class
operation. The architectural mitigation is to make expiry and rollback
explicit operations in the loop rather than retrospective hygiene.

The failure surface is real, quantified by 2024--2026 work, and largely
unmitigated in production. Every entry on it names a piece of the
optimizer that current systems do not yet have.

\subsection{11. Conclusion}\label{conclusion}

The field has built an external adaptation layer for LLMs, but not the
optimizer that should govern it. If reliability is increasingly repaired
through persistent scaffold artifacts, then production AI needs a theory
and engineering discipline for the discrete, gated, auditable learning
that happens there, not just for the differentiable learning inside the
model. This paper is one step toward that discipline: a survey of where
production sits on the \emph{patchwork to random-search to descent}
gradient of §2, a formalisation of the Stage-3 limit under stated
assumptions, and a named missing architecture where the optimizer most
plainly does not yet exist.

The field did not wait for frontier weights to become perfectly
reliable. It built a second learning system around them, made of
markdown files, skills, memories, tools, orchestration graphs, eval
suites, PRs, version histories, and rollback buttons. The mechanisms
were already in production. What changed in 2025 and 2026 is that the
combination became dense enough across vendors to read as a single
external adaptation surface, most of which is still being maintained as
patchwork. Once that surface is named, the missing pieces become
concrete engineering targets. The most consequential is governed
cross-tenant scaffold optimization: Auto-gated Loop 1 plus governed
multi-tenant cross-org Loop 2, with versioned lineage, RBAC, and
rollback. Whether it should be filled is deployment-dependent. That it
has not been is the survey's headline finding.

The frontier model generalises. The local scaffold specialises.
Reliability comes from governing that specialisation, and the most
consequential open architecture problem is the one that combines
automated local evolution with auditable cross-tenant aggregation:
turning today's patchwork into tomorrow's optimizer.

\begin{center}\rule{0.5\linewidth}{0.5pt}\end{center}

\subsection{Appendix A. Representative survey
rows}\label{appendix-a.-representative-survey-rows}

The 30 rows below are a representative spot-check of the corpus,
covering all six substrates, both research and production, and all four
clusters of §7. \emph{Loop 1} and \emph{Loop 2} columns apply to systems
satisfying the composite-two-loop criterion; \emph{n/a} means the loop
is absent. \emph{PP} is the integrated patch-local-adaptation score
(0--5); \emph{Evidence} abbreviates the strongest available source tier
(T1 peer-reviewed; T2 arXiv; T3 production claim; T4 open-source
artifact; T5 industry blog).

\begingroup\small

{\def\LTcaptype{none} % do not increment counter
\begin{longtable}[]{@{}
  >{\raggedright\arraybackslash}p{(\linewidth - 14\tabcolsep) * \real{0.1212}}
  >{\raggedleft\arraybackslash}p{(\linewidth - 14\tabcolsep) * \real{0.0404}}
  >{\raggedright\arraybackslash}p{(\linewidth - 14\tabcolsep) * \real{0.1111}}
  >{\raggedright\arraybackslash}p{(\linewidth - 14\tabcolsep) * \real{0.0707}}
  >{\raggedright\arraybackslash}p{(\linewidth - 14\tabcolsep) * \real{0.0707}}
  >{\raggedleft\arraybackslash}p{(\linewidth - 14\tabcolsep) * \real{0.0404}}
  >{\raggedright\arraybackslash}p{(\linewidth - 14\tabcolsep) * \real{0.1818}}
  >{\raggedright\arraybackslash}p{(\linewidth - 14\tabcolsep) * \real{0.3636}}@{}}
\toprule\noalign{}
\begin{minipage}[b]{\linewidth}\raggedright
System
\end{minipage} & \begin{minipage}[b]{\linewidth}\raggedleft
Year
\end{minipage} & \begin{minipage}[b]{\linewidth}\raggedright
Substrate
\end{minipage} & \begin{minipage}[b]{\linewidth}\raggedright
Loop 1
\end{minipage} & \begin{minipage}[b]{\linewidth}\raggedright
Loop 2
\end{minipage} & \begin{minipage}[b]{\linewidth}\raggedleft
PP
\end{minipage} & \begin{minipage}[b]{\linewidth}\raggedright
Evidence
\end{minipage} & \begin{minipage}[b]{\linewidth}\raggedright
Why included
\end{minipage} \\
\midrule\noalign{}
\endhead
\bottomrule\noalign{}
\endlastfoot
NanoResearch & 2026 & Integrated & Auto & Auto & 5 & T2
(arXiv:2605.10813) & Tri-level Skills/Memory/Policy co-evolution;
cleanest full-auto exemplar \\
AutoAgent & 2026 & Integrated & Auto & Auto & 5 & T2 (arXiv:2603.09716)
& Dual-cycle Execution+Evolution + elastic memory orchestration \\
SkillRL & 2026 & Integrated & Auto & Auto & 5 & T2 (arXiv:2602.08234) &
Recursive skill-augmented RL with \(\text{SR}<0.4\) threshold gate \\
AlphaEvolve & 2025 & Integrated & Auto & Auto & 4 & T2/T3
(arXiv:2506.13131) & Production proof-point: Borg +0.7\%, Gemini kernel
+23\%, FlashAttention +32.5\% \\
DGM & 2025 & Integrated & Auto & Auto & 4 & T2 (arXiv:2505.22954) &
Recursive self-modifying code; SWE-bench 20\%→50\% \\
ADAS & 2024 & Integrated & Auto & Auto & 4 & T2 (arXiv:2408.08435) &
Meta-agent search; Turing-complete agent representation \\
SkillForge & 2026 & Integrated & Auto & Human & 4 & T2
(arXiv:2604.08618) & Production hybrid exemplar; LLM-judge
\textgreater90\% + VFS versioning \\
CASCADE & 2025 & Integrated & Auto & Human & 4.5 & T2 (arXiv:2512.23880)
& Largest within-benchmark ablation gain: +57.9 pp on SciSkillBench \\
EvolveR & 2025 & Integrated & Auto & Auto & 3 & T2 (arXiv:2510.16079) &
Experience-driven lifecycle; accepted at ICML 2026 \\
Voyager & 2023 & Integrated & Auto & Auto & 4 & T2 (arXiv:2305.16291) &
Foundational skill-library result; 3.3× items, 15.3× tech-tree
milestones \\
AGENTS.md & 2025 & S1 Instructions & n/a & n/a & 2 & T3/T4 (agents.md) &
Cross-vendor standard; 60k+ repos; Lulla 2026 measures −28.6\%
runtime \\
Claude Code CLAUDE.md & 2025 & S1 Instructions & n/a & n/a & 4 & T3/T4
(docs.anthropic.com) & Four scoping levels including org-IT policy +
auto-memory layer; ceiling for S1 \\
Cursor Rules & 2024 & S1 Instructions & n/a & n/a & 3 & T3/T4
(\url{docs.cursor.com/rules}) & \texttt{.cursor/rules/*.mdc} multi-scope
+ MCP attach; git-tracked \\
Anthropic Agent Skills & 2025 & S2 Skills & n/a & Human & 3 & T3
(\url{anthropic.com/engineering}) & Progressive disclosure spec;
LangChain replication 29\%→95\% pass-rate \\
SAGE & 2025 & S2 Skills & Auto & Auto & 5 & T2 (arXiv:2512.17102) &
+8.9\% SGC, 26\% fewer steps, 59\% fewer tokens on AppWorld \\
COSPLAY & 2026 & S2 Skills & Auto & Auto & 5 & T2 (arXiv:2604.20987) &
Co-evolving decision agent + learnable skill bank; long-horizon tasks \\
Reflexion & 2023 & S3 Memory & Auto & n/a & 4 & T1 (NeurIPS 2023;
arXiv:2303.11366) & Canonical failure-triggered episodic memory; +8\%
HotpotQA \\
Generative Agents & 2023 & S3 Memory & Auto & n/a & 4 & T1 (UIST 2023;
arXiv:2304.03442) & Reflection + importance memory primitive \\
Cuadros et al.~governed collaborative memory & 2026 & S3 Memory & Auto &
Auto & 5 & T2 (arXiv:2605.04264) & Only surveyed memory system with full
provenance + versioning + correction \\
MCP & 2024 & S4 Tools & n/a & n/a & 5 & T3 (modelcontextprotocol.io) &
Cross-vendor standard; thousands of public servers; substantial
enterprise adoption \\
Harvey & 2026 & S4 Tools & n/a & n/a & 5 & T3 (harvey.ai) & 400K
queries/day; 18,000+ workflows; 200+ legal data sources \\
Hippocratic AI & 2026 & S4 Tools & n/a & n/a & 5 & T3
(hippocraticai.com) & \$3.5B valuation; 30\% readmission reduction;
360\% care-capacity boost \\
Multi-Agent Evolve (MAE) & 2025 & S5 Orchestration & Auto & n/a & 4 & T2
(arXiv:2510.23595) & RL co-evolution of Proposer/Solver/Judge
population; one loop closed, no Loop-2 \\
Cemri et al.~(AG2) & 2025 & S5 Orchestration & n/a & n/a & 3 & T2
(arXiv:2503.13657) & Specialisation +4.5 pp on GSM-Plus (\(p = 0.03\));
MAST 14 failure modes \\
Braintrust & 2024 & S6 Governance & n/a & Auto & 5 & T3 (braintrust.dev)
& Prompts-as-code: immutable commits, eval-gated PR merge, sub-5 min
rollback \\
Vellum & 2024 & S6 Governance & n/a & Auto & 5 & T3 (vellum.ai) & Same
governance pattern; production deployment with eval thresholds \\
LangSmith Hub & 2024 & S6 Governance & n/a & Auto & 5 & T3
(\url{smith.langchain.com}) & Prompt repository with environment
promotion and eval blocking \\
Self-Refine* & 2023 & (contrast) & n/a & n/a & 1 & T2 (arXiv:2303.17651)
& Fails the patch-local discriminator: in-context iteration only \\
Mem0* & 2024 & (contrast) & n/a & n/a & 2 & T3/T4 (mem0.ai) & Fails Loop
2: per-user memory with no cross-context promotion path \\
Tran \& Kiela* & 2026 & (contrast) & n/a & n/a & 1 & T2
(arXiv:2604.02460) & Single-agent baselines match or beat multi-agent
under matched thinking-token budgets; orthogonal finding constraining S5
claims \\
\end{longtable}
}

\endgroup

\emph{Asterisks mark contrast systems included for definitional
clarity.}

\begin{center}\rule{0.5\linewidth}{0.5pt}\end{center}

\subsection{Appendix B. Data appendix}\label{appendix-b.-data-appendix}

\textbf{Corpus size.} The survey scores approximately 130 unique LLM
systems across the six substrates. Each system may be scored on multiple
substrates; the per-substrate row count therefore exceeds the
unique-system count.

\textbf{Core / contrast split.} Approximately 90 unique systems satisfy
the patch-local discriminator (score ≥ 3 on at least one substrate) and
form the core corpus. The remaining \textasciitilde38 systems form the
contrast corpus: commonly described as agentic or self-improving in
public discourse but failing at least one discriminator. Appendix D
enumerates representative contrast systems.

\textbf{Inclusion criteria.} Any system that (a) wraps a foundation
model with at least one persistent artifact updated from deployment
signal, or (b) is named in the public discourse as a ``self-improving,''
``scaffolded,'' or ``agentic'' system. The contrast set (Self-Refine,
ReAct, plain Mem0, DSPy/TextGrad) is included so the discriminator has
something to mark against.

\textbf{Exclusion criteria.} Systems whose only update mechanism is
weight-level fine-tuning or RLHF: these are model-side, not
scaffold-side, and belong to the self-evolution literature this paper
differentiates from. Systems with no public evidence of deployment
(white papers without code, demos, or production claim).

\textbf{Evidence tiers.} Each system carries one or more of: (T1)
peer-reviewed publication; (T2) arXiv preprint; (T3) production claim
with metrics; (T4) open-source artifact; (T5) industry-blog or
interview-level documentation. T1/T4 carry highest weight in cluster
analysis; T3 is admissible where the claim is concrete (specific
benchmark, named integration target); T5 is admissible only when
corroborated by T1--T4 elsewhere. Tier flags are propagated to cluster
claims through the per-row evidence column in Appendix A; T3 production
claims are not pooled with T1/T2 evidence when computing cluster-level
statistics.

\textbf{Evidence-tier distribution.} Of the core corpus, approximately
one quarter carry T1 evidence, roughly half carry T2 evidence, and the
remainder carry T3 or T4 evidence as the primary tier.

\textbf{Survey window and dating.} January 2023 -- May 2026.
Preprint-only systems carry T2 evidence and are flagged in Appendix A. A
fraction of T2 work is recent enough that ablations may yet be revised;
cluster-level claims are robust to dropping any single T2 system.

\textbf{Scoring discipline.} Each system received a
substrate-by-substrate 0--5 score against the rubric of §4 and an
integrated patch-local-adaptation score (PP) in \{0,1,2,3,3.5,4,4.5,5\}.
Half-scores were used sparingly when one criterion was clearly met and
another partially. Ambiguous cases were resolved conservatively
(downward). Scoring is single-rater; Appendix A and Appendix D together
expose representative per-system scores and exclusion reasons so readers
can audit ambiguous classifications independently.

\textbf{Composite two-loop criterion (sensitivity analysis).} The strict
count of thirteen research systems satisfying the composite-two-loop
criterion changes under two natural relaxations. Relaxation (a) counts
any of versioning, lineage, or rollback as Loop-2 governance; the count
rises to approximately seventeen by re-admitting MAE, MetaGen,
MetaReflection, and CLIN. Relaxation (b) admits session signal that
updates optimizer state without persisting an artifact distinct from
optimiser state; the count rises to approximately fifteen by
re-admitting DSPy and TextGrad. The qualitative findings (Auto×Auto
well-populated; {[}Human × Human{]} and {[}Human × Auto{]} empty in the
surveyed corpus; governance-first cluster in {[}None × Auto{]}; the
cross-tenant governed-optimisation cell empty) survive both relaxations.

\begin{center}\rule{0.5\linewidth}{0.5pt}\end{center}

\subsection{Appendix C. Gate type → estimator-family
mapping}\label{appendix-c.-gate-type-estimator-family-mapping}

The §3 formalism collapses production gates into a single per-edit gate
\(G_t\) that estimates the directional loss change
\(\widehat{\Delta}_{z_t} L_D\). In practice, surveyed gates instantiate
that abstraction in several distinct ways, and some are better read as
\emph{proposal-pruning} steps that precede the per-edit gate rather than
as direct \(\widehat{\Delta}\) estimators.

{\def\LTcaptype{none} % do not increment counter
\begin{longtable}[]{@{}
  >{\raggedright\arraybackslash}p{(\linewidth - 6\tabcolsep) * \real{0.2500}}
  >{\raggedright\arraybackslash}p{(\linewidth - 6\tabcolsep) * \real{0.2500}}
  >{\raggedright\arraybackslash}p{(\linewidth - 6\tabcolsep) * \real{0.2500}}
  >{\raggedright\arraybackslash}p{(\linewidth - 6\tabcolsep) * \real{0.2500}}@{}}
\toprule\noalign{}
\begin{minipage}[b]{\linewidth}\raggedright
Gate family
\end{minipage} & \begin{minipage}[b]{\linewidth}\raggedright
What it estimates
\end{minipage} & \begin{minipage}[b]{\linewidth}\raggedright
Slot in §3 formalism
\end{minipage} & \begin{minipage}[b]{\linewidth}\raggedright
Representative systems
\end{minipage} \\
\midrule\noalign{}
\endhead
\bottomrule\noalign{}
\endlastfoot
\textbf{Eval-cascade} & \(\widehat{\Delta} L_D\) on a held-out task
suite; cascaded through cheap-to-expensive evaluators & Per-edit gate
\(G_t\) & AlphaEvolve, Braintrust, Vellum, LangSmith Hub \\
\textbf{RL-threshold} & Reward improvement vs.~a deployment threshold &
Per-edit gate \(G_t\) & SkillRL, SAGE, EvolveR \\
\textbf{LLM-judge with calibration} & Pairwise- or rubric-scored
quality, calibrated against human agreement & Per-edit gate \(G_t\) &
SkillForge (Loop 1), CASCADE (consolidation gate) \\
\textbf{Surrogate verifier} & Symbolic, learned, or test-case verifier
producing pass/fail or graded score & Per-edit gate \(G_t\) & DGM
(sandboxed code), Voyager (in-game verifier) \\
\textbf{Reviewer-judgment} & Human credit assignment via PR review or
expert sign-off & Per-edit gate \(G_t\) (human-instantiated) &
SkillForge (Loop 2), Sierra Agent OS 2.0, Anthropic skills repo \\
\textbf{Semantic-merge} & Whether a candidate is semantically duplicated
by, or compatible with, existing scaffold artifacts &
\emph{Proposal-pruning} preceding \(G_t\) & NanoResearch (orchestrator
semantic-merge step) \\
\textbf{MAP-Elites / archive selection} & Whether a candidate populates
a previously unfilled cell in a behavior-diversity archive &
\emph{Proposal-pruning} preceding \(G_t\) & AlphaEvolve (archive), ADAS
(meta-agent search) \\
\textbf{Multi-sample agreement} & Whether \(k\) independent proposals or
evaluations converge on the same direction; false-accept rate decays as
\(p^k\) rather than variance shrinking as \(1/k\) & \emph{Aggregation
filter} (qualitatively distinct from mini-batch averaging) &
Under-explored in the surveyed corpus \\
\end{longtable}
}

The eval-cascade, RL-threshold, LLM-judge, surrogate-verifier, and
reviewer-judgment families directly populate the per-edit gate \(G_t\)
of §3 and inherit its descent guarantees (under the assumptions of
Appendix E). The semantic-merge and archive-selection families are not
estimators of \(\widehat{\Delta} L_D\). They constrain the proposal
distribution rather than judging individual candidate quality, and the
descent argument therefore applies to the composition (proposal-pruning
∘ per-edit gate). Multi-sample agreement filters are listed for
completeness because they are the correct ML analogue for ``wait for
\(k\) pieces of evidence before accepting an edit,'' which is \emph{not}
the same operation as mini-batch averaging and should not be conflated
with it.

\begin{center}\rule{0.5\linewidth}{0.5pt}\end{center}

\subsection{Appendix D. Contrast set}\label{appendix-d.-contrast-set}

The contrast corpus of approximately 38 systems is included to make the
patch-local discriminator audit-able. Representative entries:

{\def\LTcaptype{none} % do not increment counter
\begin{longtable}[]{@{}
  >{\raggedright\arraybackslash}p{(\linewidth - 6\tabcolsep) * \real{0.2500}}
  >{\raggedright\arraybackslash}p{(\linewidth - 6\tabcolsep) * \real{0.2500}}
  >{\raggedright\arraybackslash}p{(\linewidth - 6\tabcolsep) * \real{0.2500}}
  >{\raggedright\arraybackslash}p{(\linewidth - 6\tabcolsep) * \real{0.2500}}@{}}
\toprule\noalign{}
\begin{minipage}[b]{\linewidth}\raggedright
System
\end{minipage} & \begin{minipage}[b]{\linewidth}\raggedright
Year
\end{minipage} & \begin{minipage}[b]{\linewidth}\raggedright
Primary discriminator failure
\end{minipage} & \begin{minipage}[b]{\linewidth}\raggedright
Brief reason
\end{minipage} \\
\midrule\noalign{}
\endhead
\bottomrule\noalign{}
\endlastfoot
Self-Refine (Madaan et al.) & 2023 & No persistent artifact & In-context
iteration only \\
ReAct (Yao et al.) & 2023 & No persistent artifact & Thought-act-observe
loop is intra-session \\
Reflexion (Shinn et al.) & 2023 & No Loop 2 & Per-agent episodic memory
without promotion \\
Generative Agents & 2023 & No Loop 2 & Per-agent reflection + importance
memory \\
ExpeL & 2023 & Loop boundary blurred & Experience bank built from
training tasks \\
Mem0 (vanilla) & 2024 & No Loop 2 & Per-user memory; no cross-context
promotion \\
Zep & 2024 & No Loop 2 & Per-agent/session knowledge graph \\
ChatGPT Memory & 2024 & No Loop 2 & Per-user persistence \\
Claude Memory & 2025 & No Loop 2 & Per-conversation-thread
persistence \\
MemGPT & 2023 & No Loop 2 & Virtual-context management per agent \\
DSPy & 2023 & No artifact distinct from optimiser state &
Prompt-template optimization \\
TextGrad & 2024 & No artifact distinct from optimiser state & Same
boundary case as DSPy \\
MetaGen & 2026 & No Loop 2; no cross-session persistence & Role pool
dynamic within single inference call \\
MAE (Multi-Agent Evolve) & 2025 & No Loop 2 & Co-evolution within
Proposer/Solver/Judge triplet \\
MetaReflection & 2024 & No Loop 2 & Per-agent semantic memory \\
CLIN & 2023 & No cross-instance promotion & Causal abstractions
per-agent-per-environment \\
LangGraph, AutoGen, CrewAI & 2023--2024 & Framework only & Orchestration
frameworks; not deployed systems \\
Tran \& Kiela & 2026 & Single-agent baseline & Constrains S5 maturity
claims (see Appendix F) \\
W\&B Weave, Agenta, LangFuse & 2024 & No Loop 1 & Observability /
eval-as-code \\
Mobile-Agent-E, Cradle, OS-Copilot & 2024--2025 & No Loop 2 &
Self-evolving but per-session/per-deployment \\
Agent S, Agent S2 & 2024--2025 & No Loop 2 & Computer-use agentic
framework \\
ACE (Zhang et al.) & 2025 & No Loop 2 in surveyed form & Single-context
evolving playbook \\
Cursor Rules (vanilla) & 2024 & No Loop 2 & Per-project artifacts \\
GitHub Copilot custom instructions, Windsurf, Continue.dev, Zed AI rules
& 2024 & No Loop 1 & Engineer-authored at repo base branch \\
Replit \texttt{replit.md} & 2025 & Loop 2 absent & Self-writing scaffold
but per-workspace \\
Aider conventions, Lovable knowledge & 2024 & No Loop 1 &
Engineer-authored convention files \\
\end{longtable}
}

The representative entries above cover the main exclusion families
(in-context-only; no Loop 2; framework-only; engineer-authored without
failure trigger; optimiser-state-only); finer-grained per-system
breakdowns are deferred to follow-on practitioner work.

\begin{center}\rule{0.5\linewidth}{0.5pt}\end{center}

\subsection{Appendix E. Formalization of artifact-layer
descent}\label{appendix-e.-formalization-of-artifact-layer-descent}

This appendix supplies the verification machinery for §3. The body
states the descent picture as a mechanism; this appendix derives the
one-step descent inequality, states the convergence proposition with its
assumptions, gives the full optimisation-analogy correspondence, walks
four concrete cases, and clarifies what credit assignment without chain
rule means.

\subsubsection{E.1 The mechanism stack}\label{e.1-the-mechanism-stack}

The abstract update rule of §3 expands into six operations that every
patch-local system must perform, whether explicitly engineered or
accidentally emergent.

\begin{enumerate}
\def\labelenumi{\arabic{enumi}.}
\tightlist
\item
  \textbf{Residual capture.} Log failures, corrections, rejected
  outputs, user edits, eval regressions, tool-call failures. The Loop-1
  signal source. Without it, \(z_t\) has no provenance.
\item
  \textbf{Mode discovery.} Cluster repeated failures into patch-specific
  error modes: the \emph{Architecture of Errors} catalogue, instantiated
  locally per patch.
\item
  \textbf{Delta synthesis.} Convert a recurring mode into a candidate
  scaffold edit: an instruction rewrite, a new skill, a memory
  correction, a tool binding, an orchestration change, an eval case.
  This produces \(z_t\).
\item
  \textbf{Delta validation.} Estimate whether the edit reduces \(L_D\)
  without collateral damage on adjacent inputs. The gate \(G_t\) is
  Auto-gated when the check is an algorithmic criterion, Human-gated
  when it is reviewer judgment.
\item
  \textbf{Promotion control.} Decide whether the accepted edit stays in
  the local fork or climbs the scope hierarchy: project, stack,
  organisation, cross-tenant. Loop 2.
\item
  \textbf{Drift and bloat control.} Expire stale rules, prune
  conflicting memories, roll back regressed skills, detect over-fitting.
  Without this, accumulated edits produce the scaffold-level analogue of
  plasticity loss.
\end{enumerate}

Each operation is the locus of a separate engineering discipline.
Surveyed systems implement them partially and unevenly: research
full-auto systems excel at (3) and (4); production governance leaders at
(5); industry hybrids at most of (1)--(5) but rarely (6). Implementation
patterns are deferred to follow-on practitioner work.

\subsubsection{E.2 Setup and notation}\label{e.2-setup-and-notation}

Let the frontier model have fixed parameters \(\theta_M\). Let the
scaffold be a structured parameter object
\[\theta_s \in \Theta_s = \Theta_1 \times \Theta_2 \times \Theta_3 \times \Theta_4 \times \Theta_5 \times \Theta_6,\]
where the six coordinates correspond to the six substrates of §4. The
space \(\Theta_s\) is discrete, symbolic, and versioned. For a
deployment patch \(D\), define patch loss as
\[L_D(\theta_s) = \mathbb{E}_{x \sim D}\!\left[\ell(\pi(\theta_M, \theta_s, x))\right].\]

A candidate edit \(z_t \in \mathcal{A}(\theta_s^t)\) is drawn from a
proposal process \(z_t \sim Q_t(z \mid H_t, \theta_s^t)\), where \(H_t\)
is the accumulated history of failures, corrections, rejected outputs,
eval regressions, reviewer comments, tool-call traces, and prior
accepted deltas, and conditioning on \(\theta_s^t\) ensures
admissibility at the current scaffold. The finite directional difference
induced by the edit is
\[\Delta_z L_D(\theta_s) = L_D(\theta_s \oplus z) - L_D(\theta_s).\] A
useful edit is one with \(\Delta_z L_D < 0\). The gate's noisy estimate
of this quantity is \(\widehat{\Delta}_z L_D\).

\subsubsection{E.3 Descent inequality}\label{e.3-descent-inequality}

The gate operates on the \emph{estimator}, not the truth. The descent
argument distinguishes them.

\textbf{Assumption A1 (positive descent bias on accepted edits).} There
exists \(\gamma_t > 0\) such that
\[\mathbb{E}\!\left[\widehat{\Delta}_{z_t} L_D \mid G_t = 1, \theta_s^t\right] \leq -\gamma_t - \tau_t.\]
The gate's accept-threshold \(\tau_t\) ensures accepted edits clear the
margin under the estimator.

\textbf{Assumption A2 (bounded estimator bias on accepted edits).} There
exists \(\varepsilon_t \geq 0\) such that
\[\big|\mathbb{E}\!\left[\Delta_{z_t} L_D - \widehat{\Delta}_{z_t} L_D \mid G_t = 1, \theta_s^t\right]\big| \leq \varepsilon_t.\]

Combining A1 and A2 with the gated update rule of §3 and writing
\(p_t = \Pr(G_t = 1 \mid \theta_s^t)\),
\[\mathbb{E}\!\left[L_D(\theta_s^{t+1}) \mid \theta_s^t\right] \leq L_D(\theta_s^t) - p_t(\gamma_t + \tau_t) + p_t \varepsilon_t.\]
This is descent under a noisy estimator: per-step expected improvement
is the gate's margin minus its estimator bias, weighted by acceptance
probability.

\subsubsection{E.4 Convergence
proposition}\label{e.4-convergence-proposition}

\textbf{Proposition (convergence of loss, informal).} Assume A1, A2,
bounded loss \(0 \leq L_D \leq B\), and additionally:

\begin{itemize}
\tightlist
\item
  \textbf{A3 (proposal coverage).} At every non-\(G\)-stable scaffold
  there exists \(\delta > 0\) such that \(Q_t\) places probability at
  least \(q > 0\) on an admissible edit with
  \(\mathbb{E}[\Delta_z L_D] \leq -\delta\).
\item
  \textbf{A4 (effective compactness).} The reachable scaffold orbit
  takes values in a finite or precompact subset of \(\Theta_s\).
\item
  \textbf{A5 (persistent acceptance with summable bias).}
  \(\sum_t p_t (\gamma_t + \tau_t) = \infty\) and
  \(\sum_t p_t \varepsilon_t < \infty\).
\end{itemize}

Then \(L_D(\theta_s^t)\) is a non-negative supermartingale up to
summable noise; \(L_D(\theta_s^t) \to L_\infty\) almost surely for some
random \(L_\infty \geq 0\); and the limit set of \(\{\theta_s^t\}\) is
\textbf{\(G\)-stable}: no admissible edit clears the gate with negative
expected directional loss in the limit.

The proposition claims convergence of the \emph{loss} and stability of
the \emph{gate decision}, not convergence to a global or local optimum
of \(L_D\). A globally suboptimal \(G\)-stable scaffold is consistent
with the result: \(Q_t\) may never propose the improving edit, or
\(\widehat{\Delta}\) may be biased to reject an improving direction.
This is the ordinary target of local stochastic optimization under noisy
estimators rather than a guarantee of optimality.

\subsubsection{E.5 Optimisation-analogy
correspondence}\label{e.5-optimisation-analogy-correspondence}

{\def\LTcaptype{none} % do not increment counter
\begin{longtable}[]{@{}
  >{\raggedright\arraybackslash}p{(\linewidth - 2\tabcolsep) * \real{0.5000}}
  >{\raggedright\arraybackslash}p{(\linewidth - 2\tabcolsep) * \real{0.5000}}@{}}
\toprule\noalign{}
\begin{minipage}[b]{\linewidth}\raggedright
Optimization concept
\end{minipage} & \begin{minipage}[b]{\linewidth}\raggedright
Scaffold analogue
\end{minipage} \\
\midrule\noalign{}
\endhead
\bottomrule\noalign{}
\endlastfoot
Parameter vector & Instructions, skills, memories, tools, routing
graphs, eval gates, governance rules \\
Forward pass & Model--scaffold composite executing on patch input \\
Loss & Failure, user correction, eval regression, retrieval miss, tool
error, schema violation \\
Finite directional difference &
\(\Delta_z L_D(\theta_s) = L_D(\theta_s \oplus z) - L_D(\theta_s)\) \\
Gradient estimate & Credit assignment identifying the scaffold
coordinate responsible for the failure \\
SGD step & Accepted scaffold edit \\
Aggregated update & Multiple recurring failures aggregated before update
(either by sampling, \emph{mini-batch averaging}, or by
sign-of-direction agreement across \(k\) proposals, \emph{multi-sample
voting}) \\
Learning rate & Edit magnitude under \(\oplus\) \\
Update radius & Promotion scope: which loss is being descended on (USER,
PROJECT, STACK, TENANT, CORE) \\
Validation loss & Held-out patch eval before merge \\
Regularization & Complexity penalty \(\Omega(z)\), bloat control,
cross-patch regression checks \\
Overfitting & Scaffold improves on this patch and degrades elsewhere \\
Rollback & Reverting a harmful accepted update \\
\end{longtable}
}

\subsubsection{E.6 Credit assignment without chain
rule}\label{e.6-credit-assignment-without-chain-rule}

Scaffold systems do not compute chain-rule derivatives through markdown
files or tool registries, and the correspondence above deliberately
omits a ``backpropagation'' row. They do, however, perform \emph{credit
assignment} through a pipeline. If the final answer fails, the
responsible coordinate may be a retrieval configuration three layers
upstream, a missing tool schema constraint, an underspecified
instruction, a stale memory, or a bad orchestration edge. Engineers do
this manually; LLM-judges and surrogate verifiers do it
semi-automatically. The operation is \emph{attribution} (output error is
assigned to an upstream coordinate) rather than backpropagation. The
descent argument of §3 needs only that this attribution succeeds with
non-zero probability per failure (assumption (iv) of the scope box), not
that it is performed by chain rule.

\subsubsection{E.7 Worked cases}\label{e.7-worked-cases}

\textbf{RAG.} Suppose failures recur because retrieved context omits the
decisive clause. The observed loss is not ``the model hallucinated.''
Credit assignment traces the error to the retrieval coordinate:
chunking, query rewrite, metadata filter, embedding model, reranker,
citation policy, or context-packing rule. The candidate edit might
tighten chunk boundaries, add a metadata constraint, change the
reranker, or insert a regression eval containing the missed clause. The
gate tests whether retrieval recall and answer accuracy improve on the
patch eval without degrading adjacent cases. If accepted, the scaffold
has moved downhill on the local loss surface.

\textbf{Pitch review.} Suppose the system repeatedly overrates decks
claiming a ``huge market'' without numeric evidence. The failure cluster
identifies an under-specified rubric coordinate. The candidate edit
changes ``evaluate TAM size'' into ``require a numeric market estimate
and source; if absent, cap TAM score at 4/10.'' The gate runs the
revised rubric against held-out decks with human scores. If agreement
improves without damaging unrelated criteria, the edit is committed.

\textbf{Tool-using agent.} Suppose calls repeatedly fail because
optional parameters are underspecified. The responsible coordinate may
be the tool schema, the call policy, or the clarification rule. The
candidate edit adds required fields, constrained decoding, or a pre-call
uncertainty check. The gate estimates whether tool-call failure
decreases without increasing refusal, latency, or unnecessary
clarification.

\textbf{Memory.} Suppose a stale customer preference is repeatedly
retrieved and causes bad decisions. The failure is not a weight error.
The responsible coordinate is a memory entry plus its provenance,
priority, and expiry policy. The candidate edit corrects, expires, or
version-tags the memory. The gate checks whether the correction improves
future behavior without erasing useful history.

\subsubsection{E.8 Loop 1 and Loop 2 as descent steps of different
radius}\label{e.8-loop-1-and-loop-2-as-descent-steps-of-different-radius}

Loop 1 proposes and validates local descent steps on a single \(L_D\).
Loop 2 expands the update radius. It selects, merges, and promotes
accepted local edits so that they apply to a larger averaging set of
patch losses. USER-level acceptance descends on a single \(L_D\).
PROJECT promotion descends on a project-averaged loss. STACK promotion
widens the average further. CORE or cross-tenant promotion descends on a
federated average. The two-loop design space is therefore not just a
taxonomy of engineering patterns: it is a taxonomy of \emph{which loss
each surveyed system is descending on}, gated by which estimator.

\subsubsection{E.9 Overfitting in scaffold
adaptation}\label{e.9-overfitting-in-scaffold-adaptation}

In weight training, overfitting is usually a defect. In scaffold
adaptation, local overfitting is partly the point. Frontier weights are
trained to preserve cross-patch generality; they must smooth over local
conventions and rare workflow-specific failures. A deployment patch
needs the opposite: selective specialization to local residuals. The
scaffold is therefore an intentionally overfittable layer, but one whose
overfitting is scoped, inspectable, versioned, gated, and reversible. A
scaffold that improves one repository, hospital workflow, or legal
diligence process may degrade performance elsewhere. That is not a
contradiction; it is patch adaptation. The mistake is promoting that
edit beyond its evidence radius.

In non-stationary patches (\(D = D_t\)), the same descent mechanism
tracks a moving optimum rather than converging once. This is why recency
weighting, replay, pruning, expiry, and rollback are not optional
hygiene; they are the scaffold-level analogues of continual-learning
machinery.

\begin{center}\rule{0.5\linewidth}{0.5pt}\end{center}

\subsection{Appendix F. Survey detail}\label{appendix-f.-survey-detail}

This appendix carries the per-substrate distributions, the per-cluster
system walks, the cell-by-cell two-loop matrix discussion, the
closest-approaches enumeration for the missing cross-tenant
governed-optimisation architecture, and the orchestration-bounding
result of Tran \& Kiela. It supplies the verification material for
§§6--9.

\subsubsection{F.1 Substrate maturity
distribution}\label{f.1-substrate-maturity-distribution}

{\def\LTcaptype{none} % do not increment counter
\begin{longtable}[]{@{}
  >{\raggedright\arraybackslash}p{(\linewidth - 8\tabcolsep) * \real{0.1400}}
  >{\raggedleft\arraybackslash}p{(\linewidth - 8\tabcolsep) * \real{0.0400}}
  >{\raggedleft\arraybackslash}p{(\linewidth - 8\tabcolsep) * \real{0.0500}}
  >{\raggedright\arraybackslash}p{(\linewidth - 8\tabcolsep) * \real{0.2200}}
  >{\raggedright\arraybackslash}p{(\linewidth - 8\tabcolsep) * \real{0.5500}}@{}}
\toprule\noalign{}
\begin{minipage}[b]{\linewidth}\raggedright
Substrate
\end{minipage} & \begin{minipage}[b]{\linewidth}\raggedleft
\(n\)
\end{minipage} & \begin{minipage}[b]{\linewidth}\raggedleft
Mean
\end{minipage} & \begin{minipage}[b]{\linewidth}\raggedright
Score-5 systems
\end{minipage} & \begin{minipage}[b]{\linewidth}\raggedright
Survey finding
\end{minipage} \\
\midrule\noalign{}
\endhead
\bottomrule\noalign{}
\endlastfoot
\textbf{S1. Instructions} & 12 & 2.83 & \emph{(none)} & Rules persist
across IDEs and agents (Claude Code's CLAUDE.md, Cursor Rules,
AGENTS.md, Windsurf, Continue.dev) but no surveyed instruction system
closes both an automated inner loop and a governed two-loop promotion.
Ceiling held at 4 by Claude Code and Replit \\
\textbf{S2. Skills} & 16 & 2.62 & SAGE, COSPLAY & Research frontier is
active (failure-triggered skill update); production governance is
weaker. Anthropic's skills repository is one-pool with PR review but no
automated eval gate \\
\textbf{S3. Memory} & 19 & 3.16 & Cuadros et al.~governed collaborative
memory & Cross-session memory is widespread; \emph{correctable,
versioned} memory is rare. Production memory systems (ChatGPT Memory,
Claude Memory, Mem0, Zep, MemGPT) treat memory as append-only \\
\textbf{S4. Tools} & 19 & 3.63 & MCP, Harvey, Hippocratic AI & Most
mature production substrate. Tool bundles are the commercial unit; MCP
is the cross-vendor standard with broad public-server ecosystem and
substantial enterprise adoption \\
\textbf{S5. Orchestration} & 19 & 2.68 & \emph{(none)} & Multi-agent
topologies exist (LangGraph, AutoGen, CrewAI) but topology
\emph{evolution} is rare; MAE comes closest with RL co-evolution of
Proposer/Solver/Judge population but lacks Loop-2 cross-context
promotion and so does not satisfy the score-5 rubric \\
\textbf{S6. Governance} & 21 & 3.05 & Braintrust, Vellum, LangSmith Hub,
AGENTS.md/AAIF & High production maturity for prompts-as-code: immutable
commits, eval-gated PR merge, sub-five-minute rollback, typically
without inner-loop adaptation \\
\textbf{INTEGRATED} & 34 & 3.74 & NanoResearch, AutoAgent, SkillRL & The
high-water mark when present, but the cluster is dominated by research
systems lacking enterprise governance \\
\end{longtable}
}

\subsubsection{F.2 Cluster table with representative
systems}\label{f.2-cluster-table-with-representative-systems}

{\def\LTcaptype{none} % do not increment counter
\begin{longtable}[]{@{}
  >{\raggedright\arraybackslash}p{(\linewidth - 6\tabcolsep) * \real{0.1400}}
  >{\raggedright\arraybackslash}p{(\linewidth - 6\tabcolsep) * \real{0.3000}}
  >{\raggedright\arraybackslash}p{(\linewidth - 6\tabcolsep) * \real{0.1400}}
  >{\raggedright\arraybackslash}p{(\linewidth - 6\tabcolsep) * \real{0.4200}}@{}}
\toprule\noalign{}
\begin{minipage}[b]{\linewidth}\raggedright
Cluster
\end{minipage} & \begin{minipage}[b]{\linewidth}\raggedright
Representative systems (year, PP score)
\end{minipage} & \begin{minipage}[b]{\linewidth}\raggedright
Loop 1 / Loop 2
\end{minipage} & \begin{minipage}[b]{\linewidth}\raggedright
What the cluster proves
\end{minipage} \\
\midrule\noalign{}
\endhead
\bottomrule\noalign{}
\endlastfoot
Research full-auto & NanoResearch (2026, 5), SkillRL (2026, 5),
AlphaEvolve (2025, 4), DGM (2025, 4), ADAS (2024, 4), Voyager (2023, 4)
& Auto / Auto & Both loops can be fully automated under algorithmic
gates; benchmark gains compound; enterprise governance is uniformly
absent \\
Production hybrid & SkillForge (2026, 4), CASCADE (2025, 4.5), Sierra OS
2.0 (2025, 3)\footnote{Sierra Agent OS 2.0 also appears in the
  \emph{Open standards / vertical-bundle} discussion in F.4 because its
  commercial product wraps a vertical workflow around the standardised
  agent-OS surface; we place it primarily in \emph{Production hybrid} on
  the basis of its Loop-1/Loop-2 gate types.}, Cognition Devin (2024, 3)
& Auto / Human & Algorithmic Loop-1 gate plus reviewer-judgment Loop-2
gate is the production-viable cell; SkillForge is the architectural
exemplar \\
Governance-first & Braintrust, Vellum, LangSmith Hub (all 2024, S6 = 5)
& None / Auto & Eval-gated promotion exists; candidate generation is
manual; no patch-local inner loop \\
Open standards & AGENTS.md (2025), MCP (2024), Claude Skills (2025),
Cursor Rules (2024) & Mixed / Mixed & Vendor-cross convergence on
artifact format and tool attachment; tens of thousands of AGENTS.md
repositories; broad MCP server ecosystem with substantial enterprise
adoption \\
\end{longtable}
}

\subsubsection{F.3 Surprising absences}\label{f.3-surprising-absences}

Five absences stand out across the patch-local core corpus:

\begin{enumerate}
\def\labelenumi{\arabic{enumi}.}
\tightlist
\item
  \textbf{Failure-triggered memory is rare in production.} Deployed
  memory systems prefer recording preferences and successful facts; only
  Reflexion, the Voyager skill library, Gemini Code Assist's
  PR-rejection memory, and the classical SOAR chunking pattern
  explicitly record residuals. The most obvious Loop-1 signal is the one
  production memory systems mostly ignore.
\item
  \textbf{Cross-customer skill promotion is absent.} AGENTS.md is
  cross-vendor; Anthropic's skills repository is one-pool. No surveyed
  system has a vetted cross-organisational marketplace with eval-gated
  promotion.
\item
  \textbf{Memory versioning is mostly unsolved.} Only the Cuadros et
  al.~(2026) governed collaborative memory framework provides explicit
  provenance, versioning, and correction pathways.
\item
  \textbf{Scaffold-level curriculum is absent.} No surveyed system
  selects \emph{which} sessions inform the gradient. Every session
  contributes equally; weighting by recency, severity, or
  representativeness has not been published.
\item
  \textbf{Overfitting detection at Loop-2 promotion is absent.} No
  surveyed system runs a held-out eval set per project to catch
  scaffolds that over-fit one deployment before they propagate to
  another.
\end{enumerate}

\subsubsection{F.4 Per-cluster walks}\label{f.4-per-cluster-walks}

\textbf{Research full-auto.} \textbf{NanoResearch} (Xu et al., 2026)
co-evolves Skills, Memory, and Policy with a semantic-merge
orchestrator; Innovation 4.96 → 5.65 and Compliance 6.66 → 8.96 across
three rounds. \textbf{SkillRL} (arXiv:2602.08234) uses RL reward
distillation with a hard threshold at task-category accuracy below 0.4
and reports +15.3\% over strong baselines. \textbf{AlphaEvolve}
(DeepMind, 2025) runs a Gemini-driven ensemble with a cascade evaluator
and a MAP-Elites archive, deployed at Google to gain +0.7\% Borg compute
worldwide, +23\% on the Gemini training kernel, +32.5\% on
FlashAttention, and the first general-field, recursively applicable
improvement over Strassen's 1969 bound on 4×4 complex matrix
multiplication (48 scalar multiplications). The honest framing matters:
AlphaTensor (Fawzi et al., 2022) had already found a 47-multiplication
algorithm over GF(2), and Winograd reported a 48-multiplication method
for commutative rings that was not recursively applicable; AlphaEvolve's
contribution is the first improvement that holds in general fields and
admits recursive composition. \textbf{DGM} (Zhang et al., 2025) modifies
its own source code with archive-based parent selection; SWE-bench 20\%
→ 50\%. The cluster maximises velocity. Enterprise governance is
uniformly absent.

\textbf{Production hybrid.} \textbf{SkillForge} (Alibaba Cloud,
arXiv:2604.08618) drives auto-generated skill diffs through an LLM-judge
with \textgreater90\% human agreement (the Loop-1 gate); human support
engineers review what the LLM-judge passes (the Loop-2 gate). VFS
version commits provide full lineage. \textbf{CASCADE} (Huang et al.,
2025) reports the largest within-benchmark ablation gain in the surveyed
corpus (93.3\% vs.~35.4\% on SciSkillBench, +57.9 pp), driven by memory
consolidation gated by human-agent collaboration. The headline gain is
not directly comparable to AlphaEvolve's production +0.7\% Borg or DGM's
+30 pp on the broader SWE-bench: it is measured on a single benchmark
designed to reward the very mechanism CASCADE adds, and should be read
as evidence of mechanism-on-benchmark fit rather than as a cross-cluster
effect size. \textbf{Sierra Agent OS 2.0} routes AI-driven Insights
through human-authored Expert Answers in GitHub-style Workspaces.
\textbf{Cognition Devin} ships dynamic in-session tool creation with a
67\% PR-merge rate (the merged PRs are the Loop-2 promotion). This is
the production-viable cell.

\textbf{Governance-first.} \textbf{Braintrust}, \textbf{Vellum},
\textbf{LangSmith Hub}, \textbf{W\&B Weave}, \textbf{Agenta},
\textbf{LangFuse}. The pattern is \emph{prompts-as-code}: immutable
commits, semantic-version tag pointers, PR-blocked merges on eval
regression, sub-five-minute rollback. Engineers iterate the artifacts
manually; the eval pipeline gates the promotion. We frame this as
\textbf{MLOps for scaffolds}, distinct from MLOps for weights.

\textbf{Open standards.} \textbf{AGENTS.md} (Linux Foundation, December
2025) is now in over 60,000 GitHub repositories; Lulla et al.~(2026)
report −28.64\% wall-clock and −16.58\% tokens with no quality loss
after adding an AGENTS.md scaffold. \textbf{MCP} (Anthropic, 2024) has
become the cross-vendor standard with thousands of public servers and
substantial enterprise adoption. \textbf{Claude Agent Skills},
\textbf{Cursor Rules}, \textbf{GitHub Copilot custom instructions},
\textbf{Windsurf rules}, \textbf{Continue.dev rules}, \textbf{Zed AI
rules}, \textbf{Replit \texttt{replit.md}} all converge on the same
pattern: a git-tracked markdown artifact attached to a workspace, with
an MCP-style tool layer for actions. Convergence across vendors is the
strongest evidence the architecture is real rather than a single
vendor's local optimum.

The four clusters do not partition the corpus (many systems sit between
two clusters), but they capture the architectural variation.
\textbf{Vertical-bundle vendors} (Harvey for legal, Hippocratic AI for
healthcare, Sierra for customer experience) sit between \emph{open
standards} and \emph{production hybrid}: the bundle is the commercial
unit; the underlying model is mostly the same one a competitor would
use.

\textbf{Tran \& Kiela (2026) orchestration bound.} One result from the
orchestration literature bounds how strongly the cluster discussion can
be read as evidence for multi-agent architectures specifically. Tran \&
Kiela (2026) report that single-agent LLMs match or beat multi-agent
systems on multi-hop reasoning when thinking-token budgets are matched;
the implication is that observed maturity gaps on S5 (orchestration) may
reflect the rarity of governed topology \emph{evolution} and the
conflation of agent-count with compute-budget rather than a structural
deficit of multi-agent designs. The orchestration cluster is therefore
better read as evidence that production has converged on flat agent
topologies and on tool-mediated rather than role-mediated decomposition,
not as evidence that orchestration is the binding maturity constraint.

\subsubsection{F.5 Two-loop matrix: cell-by-cell
discussion}\label{f.5-two-loop-matrix-cell-by-cell-discussion}

\textbf{{[}Auto × Auto{]}} is most populated (eleven systems). Every
member has a machine-evaluated decision criterion on both loops.
Velocity is the cluster's defining property; the lack of
cross-organisational governance is its uniform weakness.

\textbf{{[}Auto × Human{]}} is the production-viable cell: Loop 1 fires
at machine speed under an algorithmic criterion, Loop 2 requires human
approval before propagation. The pattern reconciles automation velocity
with reviewer accountability. SkillForge's LLM-judge + reviewer pipeline
is the architectural exemplar. CASCADE, Sierra, and Devin instantiate
variants; the industry exemplars (Sierra, Devin) are marked with †
because vendor documentation rather than peer-reviewed metrics is the
primary evidence.

\textbf{{[}None × Auto-gated{]}} holds the governance-first
prompts-as-code systems (Braintrust, LangSmith Hub, Vellum). They have a
real Loop 2 (eval-gated CI blocks promotion on regression), but no Loop
1 in the patch-local sense, because the candidate artifacts are authored
by engineers at their desks rather than triggered by session feedback.
The asterisks mark this: governance without inner-loop adaptation. The
cell is informative because it shows what survives when only the outer
loop is automated.

\textbf{{[}Human × Auto{]}} is empty by engineering logic: if humans
author candidates manually, automated promotion offers little leverage.
A speculative occupant would be ``expert curation + auto-distribution''
(expert-authored skills automatically gated through a held-out eval
before shipping); no published example exists.

\textbf{{[}Human × Human{]}} is empty in the surveyed corpus. A
candidate occupant would be a reviewer-authored-and-reviewer-promoted
artifact pipeline targeting solo and small-team workflows where
automation cost exceeds review latency. No public example was found
during the May 2026 survey window.

\textbf{{[}Human × None{]}} is structurally improbable: a system
requiring human authorship on Loop 1 but with no Loop 2 collapses to
``engineer-edits-a-config'' and is not patch-local in the sense of §3.

\textbf{DGM dual-mode note.} DGM operates in both {[}Auto × Auto{]}
(primary archive-based evolution) and {[}Auto × Human{]} (sandboxed
code-modification review) modes depending on deployment configuration;
we list it once in the dominant cell.

\subsubsection{F.6 Closest approaches to governed cross-tenant scaffold
optimization}\label{f.6-closest-approaches-to-governed-cross-tenant-scaffold-optimization}

The four properties of governed cross-tenant scaffold optimization are
(a) auto-gated Loop 1, (b) multi-tenant cross-org promotion, (c)
versioned lineage with RBAC, (d) rollback. Closest approaches partition
the requirements:

\begin{itemize}
\tightlist
\item
  \textbf{AlphaEvolve} has (a) and partial (b) within a single
  organisation; no cross-org promotion mechanism.
\item
  \textbf{NanoResearch} has (a) at the research level; no enterprise
  governance layer described.
\item
  \textbf{SkillForge} has (a) and (b) within one enterprise, partial (c)
  via VFS commits; (d) is undocumented in the public report.
\item
  \textbf{Sierra Agent OS 2.0} has (b) and partial (a); the GitHub-style
  Workspaces provide some lineage but not RBAC-with-rollback in the
  strict sense.
\item
  \textbf{Vellum / Braintrust / LangSmith Hub} have (c) and (d)
  explicitly but no (a) or (b) in the strict sense; these are
  governance-first systems by construction.
\end{itemize}

No surveyed system combines (b), (c), and (d) under any gate type. The
cross-organisational governance triad is itself unfilled, independently
of whether Loop 1 is auto-gated.

\subsubsection{F.7 Deployment-dependence of the missing
architecture}\label{f.7-deployment-dependence-of-the-missing-architecture}

Whether the corner \emph{should} be filled is deployment-dependent. In
safety-critical domains (medical decision support, legal advice,
financial advice), reviewer-judgment gates may be a design feature
rather than a defect, and the human latency is the point. The narrower
survey claim is what matters: no public system currently combines
automated local scaffold evolution with governed cross-organisation
promotion, versioned lineage with RBAC, and rollback. Naming the corner
converts a vague gap into a specific four-property audit criterion that
any future system can be measured against. The shortest demonstrable
path appears to build on AlphaEvolve's cascade evaluator and MAP-Elites
archive, extending federation from within-DeepMind multi-target to
cross-customer multi-tenant. The architectural surface, not the
timeline, is what this paper claims.

\begin{center}\rule{0.5\linewidth}{0.5pt}\end{center}

\subsection{References}\label{references}

AGENTS.md. (2025). AGENTS.md: A cross-vendor open standard for agent
instructions. Linux Foundation. \url{https://agents.md}

Agashe, S., Han, J., Gan, S., Yang, J., Li, A., \& Wang, X. E. (2024).
Agent S: An open agentic framework that uses computers like a human.
\emph{arXiv preprint arXiv:2410.08164}.

Agashe, S., et al.~(2025). Agent S2: A compositional
generalist-specialist framework for computer use agents. \emph{arXiv
preprint arXiv:2504.00906}.

Anthropic. (2024). Model Context Protocol specification.
\url{https://modelcontextprotocol.io}

Arbuzov, M. L., Bei, S., Dong, Z., Kalaev, D., \& Shvets, A. A. (2025).
Beyond exponential decay: Rethinking error accumulation in large
language models. \emph{arXiv preprint arXiv:2505.24187}.

Arbuzov, M. L., Shvets, A. A., \& Bei, S. (2026). The architecture of
errors: Logarithmic mode discovery and polylogarithmic intervention
budgets for long-context LLM reliability. \emph{arXiv preprint}
(forthcoming).

Cemri, M., et al.~(2025). Why do multi-agent LLM systems fail?
\emph{arXiv preprint arXiv:2503.13657}. Introduces the MAST taxonomy.

Chen, M., Li, Y., Yang, Y., Yu, S., Lin, B., \& He, X. (2024).
AutoManual: Constructing instruction manuals by LLM agents via
interactive environmental learning. \emph{NeurIPS 2024}.
arXiv:2405.16247.

Chen, Y., et al.~(2025). Multi-Agent Evolve: LLM self-improve through
co-evolution. \emph{arXiv preprint arXiv:2510.23595}.

Cuadros, D. F., Maiga, A.-A., Meskhidze, H., \& Curtis-Trudel, A.
(2026). Governed collaborative memory as artificial selection in
LLM-based multi-agent systems. \emph{arXiv preprint arXiv:2605.04264}.

Dohare, S., et al.~(2024). Loss of plasticity in deep continual
learning. \emph{Nature}, 632(8026), 768--774.

Du, P. (2026). Memory for autonomous LLM agents: Mechanisms, evaluation,
and emerging frontiers. \emph{arXiv preprint arXiv:2603.07670}.

Fang, J., Peng, Y., Zhang, X., et al.~(2025). A comprehensive survey of
self-evolving AI agents: A new paradigm bridging foundation models and
lifelong agentic systems. \emph{arXiv preprint arXiv:2508.07407}.

Fawzi, A., Balog, M., Huang, A., et al.~(2022). Discovering faster
matrix multiplication algorithms with reinforcement learning.
\emph{Nature}, 610(7930), 47--53. (DeepMind AlphaTensor.)

Gao, H.-a., Geng, J., Hua, W., et al.~(2026). A survey of self-evolving
agents: What, when, how, and where to evolve on the path to artificial
super intelligence. \emph{Transactions on Machine Learning Research}.
arXiv:2507.21046.

Harvey. (2026). Harvey product overview. \url{https://www.harvey.ai}

Hippocratic AI. (2026). Polaris clinical outcome evidence.
\url{https://www.hippocraticai.com}

Hu, S., Lu, C., \& Clune, J. (2024). Automated design of agentic systems
(ADAS). \emph{arXiv preprint arXiv:2408.08435}.

Huang, X., et al.~(2025). CASCADE: Cumulative agentic skill creation
through autonomous development and evolution. \emph{arXiv preprint
arXiv:2512.23880}.

Jaroslawicz, D., et al.~(2025). How many instructions can LLMs follow at
once? \emph{arXiv preprint arXiv:2507.11538}. Introduces the IFScale
benchmark.

Kirk, R., Mediratta, I., Nalmpantis, C., et al.~(2024). Understanding
the effects of RLHF on LLM generalisation and diversity. \emph{ICLR
2024}. arXiv:2310.06452.

Liu, N. F., et al.~(2023). Lost in the middle: How language models use
long contexts. \emph{TACL}. arXiv:2307.03172.

Liu, X., et al.~(2026). SkillForge: Forging domain-specific,
self-evolving agent skills in cloud technical support. \emph{arXiv
preprint arXiv:2604.08618}. Accepted at ACM SIGIR 2026 Industry Track.

Lulla, J. L., et al.~(2026). On the impact of AGENTS.md files on the
efficiency of AI coding agents. \emph{arXiv preprint arXiv:2601.20404}.

Madaan, A., et al.~(2023). Self-Refine: Iterative refinement with
self-feedback. \emph{NeurIPS 2023}. arXiv:2303.17651.

Novikov, A., et al.~(2025). AlphaEvolve: A coding agent for scientific
and algorithmic discovery. \emph{arXiv preprint arXiv:2506.13131}.
(DeepMind.)

Padmakumar, V., \& He, H. (2023). Does writing with language models
reduce content diversity? \emph{arXiv preprint arXiv:2309.05196}.

Park, J. S., et al.~(2023). Generative agents: Interactive simulacra of
human behavior. \emph{UIST 2023}. arXiv:2304.03442.

Shinn, N., et al.~(2023). Reflexion: Language agents with verbal
reinforcement learning. \emph{NeurIPS 2023}. arXiv:2303.11366.

Sierra. (2025). Agent OS 2.0: From answers to memory and action.
\url{https://sierra.ai/blog/agent-os-2-0}

Singh, S., et al.~(2025). The leaderboard illusion. \emph{arXiv preprint
arXiv:2504.20879}.

Sumers, T. R., Yao, S., Narasimhan, K., \& Griffiths, T. L. (2023).
Cognitive architectures for language agents (CoALA). \emph{Transactions
on Machine Learning Research}. arXiv:2309.02427.

Tran, D., \& Kiela, D. (2026). Single-agent LLMs outperform multi-agent
systems on multi-hop reasoning under equal thinking token budgets.
\emph{arXiv preprint arXiv:2604.02460}.

Wang, G., et al.~(2023). Voyager: An open-ended embodied agent with
large language models. \emph{arXiv preprint arXiv:2305.16291}.

Wang, J., et al.~(2025). Reinforcement learning for self-improving agent
with skill library. \emph{arXiv preprint arXiv:2512.17102}. Method named
SAGE.

Wang, X., et al.~(2026). AutoAgent: Evolving cognition and elastic
memory orchestration for adaptive agents. \emph{arXiv preprint
arXiv:2603.09716}.

Wu, R., et al.~(2025). EvolveR: Self-evolving LLM agents through an
experience-driven lifecycle. \emph{arXiv preprint arXiv:2510.16079}.
Accepted at ICML 2026.

Wu, X., et al.~(2026). Co-evolving LLM decision and skill bank agents
for long-horizon tasks. \emph{arXiv preprint arXiv:2604.20987}.
Framework named COSPLAY.

Xia, P., et al.~(2026). SkillRL: Evolving agents via recursive
skill-augmented reinforcement learning. \emph{arXiv preprint
arXiv:2602.08234}.

Xu, Jinhang, et al.~(2026). NanoResearch: Co-evolving skills, memory,
and policy for personalized research automation. \emph{arXiv preprint
arXiv:2605.10813}.

Xu, Zhenlin, et al.~(2026). From storage to steering: Memory control
flow attacks on LLM agents. \emph{arXiv preprint arXiv:2603.15125}.
Introduces the MCFA attack class.

Zhang, J., et al.~(2025). Darwin Gödel machine: Open-ended evolution of
self-improving agents. \emph{arXiv preprint arXiv:2505.22954}. (Body
text refers to this as DGM.)

Zhou, Y., Shu, W., Su, Y., et al.~(2026). A comprehensive survey on
agent skills: Taxonomy, techniques, and applications. \emph{arXiv
preprint arXiv:2605.07358}.

Zou, W., et al.~(2024). PoisonedRAG: Knowledge corruption attacks to
retrieval-augmented generation of large language models. \emph{arXiv
preprint arXiv:2402.07867}. Accepted at USENIX Security 2025.

\end{document}
